\documentclass[11pt,a4paper]{article}

\usepackage[utf8]{inputenc}
\usepackage[T1]{fontenc}
\usepackage[italian]{babel}
\usepackage[margin=2.4cm]{geometry}
\usepackage{hyperref}
\usepackage{array}
\usepackage{booktabs}
\usepackage{listings}
\usepackage{xcolor}
\usepackage{enumitem}

\hypersetup{
  colorlinks=true,
  linkcolor=black,
  urlcolor=blue
}

\renewcommand{\contentsname}{Indice}
\setlength{\parindent}{0pt}
\setlength{\parskip}{0.45em}
\emergencystretch=2em
\setlist[itemize]{topsep=2pt,itemsep=2pt,leftmargin=1.2cm}
\setlist[enumerate]{topsep=2pt,itemsep=2pt,leftmargin=1.2cm}

\lstset{
  basicstyle=\ttfamily\small,
  breaklines=true,
  frame=single,
  backgroundcolor=\color{black!3}
}

\title{\textbf{Relazione tecnica sulla parte di rete e sulla business logic}\\
\large Focusly -- Progetto di Reti di Calcolatori}
\author{Samuele Caporale \\ Matricola 329096 \\ A.A. 2025/2026}
\date{}

\begin{document}

\maketitle
\tableofcontents
\bigskip

\section{Introduzione}

La presente relazione descrive il funzionamento della parte di collegamento di rete dell'applicazione \emph{Focusly}, con particolare attenzione all'integrazione tra frontend React, Firebase Authentication e Cloud Firestore. L'analisi e` stata ricostruita a partire dal codice effettivamente presente nel progetto, in particolare nei moduli \texttt{frontend/src/lib/firebase.js}, \texttt{frontend/src/services/*} e \texttt{frontend/src/state/store.js}.

L'obiettivo dell'app e` gestire task di studio, gruppi tematici e milestone, mantenendo la sincronizzazione dei dati su cloud mediante autenticazione Google e persistenza su Firestore. Dal punto di vista delle reti di calcolatori, il progetto realizza una tipica architettura client-cloud: il browser esegue la SPA React e comunica direttamente con i servizi Firebase tramite connessioni HTTPS gestite dall'SDK ufficiale.

\section{Architettura di rete dell'applicazione}

\subsection{Vista generale}

L'architettura realmente utilizzata nella versione corrente e` la seguente:

\begin{lstlisting}
Browser
  -> React 18 + Vite
  -> Firebase Authentication (Google Sign-In)
  -> Cloud Firestore
\end{lstlisting}

Il frontend costituisce quindi il nodo applicativo principale. Non e` presente, nella V1, un backend custom interposto tra client e database per le operazioni di business ordinarie. Nel repository esiste una cartella \texttt{backend/}, ma il suo ruolo attuale e` solo di scaffold tecnico: espone un endpoint di health check e un layer socket iniziale, senza essere coinvolto nel flusso dati di task, gruppi o autenticazione. Questo aspetto e` importante perche' la comunicazione di rete effettiva dell'app avviene direttamente tra browser e infrastruttura Firebase.

\subsection{Tecnologie coinvolte}

\begin{center}
\begin{tabular}{p{3.2cm} p{9.3cm}}
\toprule
Componente & Funzione principale \\
\midrule
React 18 & Rendering della SPA e gestione dell'interfaccia utente \\
Zustand & Store globale per stato applicativo, bootstrap e azioni asincrone \\
Firebase Auth & Identificazione dell'utente tramite account Google \\
Cloud Firestore & Persistenza cloud di profilo utente, task, gruppi e milestone \\
localStorage & Cache locale e sorgente della migrazione iniziale verso Firestore \\
\bottomrule
\end{tabular}
\end{center}

\subsection{Tipi di comunicazione di rete}

Le interazioni di rete sono di due tipi:

\begin{itemize}
  \item autenticazione remota con Google tramite popup, mediata da Firebase Authentication;
  \item accesso al database Firestore mediante letture, scritture e batch write effettuati dall'SDK Firebase.
\end{itemize}

Non sono invece presenti nella parte operativa corrente:

\begin{itemize}
  \item API REST proprietarie per task o gruppi;
  \item token JWT gestiti dal progetto in modo manuale;
  \item WebSocket o listener realtime Firestore per l'aggiornamento continuo dell'interfaccia.
\end{itemize}

\section{Autenticazione con Firebase}

\subsection{Configurazione}

La configurazione Firebase e` centralizzata in \texttt{frontend/src/lib/firebase.js}. Il modulo legge le variabili d'ambiente \texttt{VITE\_FIREBASE\_*}, inizializza l'app Firebase e crea tre oggetti condivisi: \texttt{auth}, \texttt{db} e \texttt{googleProvider}. La scelta di concentrare questa logica in un solo file riduce l'accoppiamento e impedisce inizializzazioni multiple.

Il provider Google viene inoltre configurato con:

\begin{itemize}
  \item lingua italiana per il flusso di login;
  \item parametro \texttt{prompt = select\_account}, utile per forzare la scelta esplicita dell'account.
\end{itemize}

\subsection{Servizio di autenticazione}

Il modulo \texttt{authService.js} espone tre primitive:

\begin{itemize}
  \item \texttt{signInWithGoogle()}, che avvia \texttt{signInWithPopup()};
  \item \texttt{signOutUser()}, che invalida la sessione lato Firebase;
  \item \texttt{subscribeToAuthChanges()}, che usa \texttt{onAuthStateChanged()} per osservare il cambio di utente autenticato.
\end{itemize}

L'applicazione non interroga continuamente il server per conoscere lo stato della sessione: si appoggia invece al listener di Firebase Auth, che notifica al frontend i cambiamenti di autenticazione. Questo approccio riduce traffico superfluo e complessita` applicativa.

\subsection{Bootstrap della sessione in React}

Nel componente radice \texttt{App.jsx}, una \texttt{useEffect} registra la subscription all'autenticazione. Quando Firebase segnala un utente:

\begin{enumerate}
  \item lo store Zustand richiama \texttt{handleAuthStateChange(user)};
  \item viene creato o aggiornato il profilo utente in Firestore;
  \item viene valutata l'eventuale migrazione dei dati locali;
  \item viene caricato lo stato dell'area di lavoro;
  \item l'interfaccia passa da schermata di bootstrap alla dashboard.
\end{enumerate}

In caso di assenza di utente autenticato, lo store torna allo stato base e l'app mostra \texttt{SignInScreen}. La business logic di ingresso e uscita dall'app non e` quindi distribuita nei componenti di UI, ma concentrata nello store globale.

\section{Persistenza cloud con Firestore}

\subsection{Modello dati}

Il modello Firestore adottato e` gerarchico e multi-tenant per utente:

\begin{lstlisting}
users/{uid}
users/{uid}/tasks/{taskId}
users/{uid}/groups/{groupId}
users/{uid}/groups/{groupId}/milestones/{milestoneId}
\end{lstlisting}

La collezione \texttt{users} contiene il profilo e i metadati di sincronizzazione. Le sottocollezioni \texttt{tasks} e \texttt{groups} rappresentano il workspace dell'utente; ogni gruppo contiene a sua volta le milestone. Questa struttura garantisce isolamento logico dei dati e si integra bene con le regole di sicurezza basate su \texttt{uid}.

\subsection{Mappatura tra oggetti React e documenti Firestore}

Il modulo \texttt{firestoreMappers.js} svolge una funzione centrale. I suoi compiti principali sono:

\begin{itemize}
  \item conversione delle date JavaScript e delle stringhe ISO in \texttt{Timestamp};
  \item conversione inversa da \texttt{Timestamp} a oggetti leggibili dal frontend;
  \item normalizzazione di campi numerici, booleani e delle date pianificate;
  \item uniformazione del formato dati rispetto ai modelli \texttt{Task} e \texttt{Group}.
\end{itemize}

Questa separazione evita che i componenti React debbano conoscere il formato interno di Firestore. In termini architetturali, il mapper costituisce un adattatore tra dominio applicativo e livello di persistenza.

\subsection{Operazioni CRUD e batch write}

Il modulo piu` rilevante per la comunicazione dati e` \texttt{firestoreService.js}. Esso costruisce riferimenti documentali tramite \texttt{doc()} e \texttt{collection()}, quindi esegue:

\begin{itemize}
  \item letture singole con \texttt{getDoc()};
  \item letture multiple con \texttt{getDocs()};
  \item query ordinate con \texttt{orderBy()};
  \item query filtrate con \texttt{where()};
  \item scritture idempotenti con \texttt{setDoc()};
  \item aggiornamenti mirati con \texttt{updateDoc()};
  \item scritture atomiche multiple con \texttt{writeBatch()}.
\end{itemize}

La scelta di usare batch write per create, update e delete permette di mantenere coerenza tra i documenti applicativi e il campo \texttt{workspaceRevision} presente nel profilo utente. Ogni modifica significativa incrementa infatti tale revisione mediante \texttt{increment(1)}. In questo modo l'app puo distinguere se la cache locale e` allineata con lo stato remoto.

\subsection{Regole di sicurezza}

Le regole in \texttt{frontend/firestore.rules} definiscono una politica semplice ma corretta:

\begin{itemize}
  \item lettura e scrittura consentite solo se \texttt{request.auth != null};
  \item l'utente puo accedere esclusivamente ai documenti associati al proprio \texttt{uid};
  \item la cancellazione del profilo utente \texttt{users/\{uid\}} e` disabilitata;
  \item ogni altro percorso Firestore e` negato di default.
\end{itemize}

Dal punto di vista della sicurezza di rete, la decisione fondamentale e` spostare il controllo di accesso dal frontend alle regole server-side di Firestore. Anche se il client e` modificabile dall'utente finale, il database continua a imporre isolamento tra utenti.

\section{Business logic React e organizzazione software}

\subsection{Ruolo dello store Zustand}

Lo store definito in \texttt{state/store.js} e` il vero orchestratore dell'applicazione. Ospita sia lo stato condiviso sia le principali azioni asincrone. I campi piu` importanti sono:

\begin{itemize}
  \item \texttt{tasks}, \texttt{groups} e \texttt{selectedGroupId};
  \item \texttt{authUser}, \texttt{authStatus} e \texttt{dataStatus};
  \item \texttt{migrationStatus}, \texttt{migrationSource} e \texttt{appError};
  \item \texttt{workspaceRevision}, usato per la coerenza cache/cloud.
\end{itemize}

La logica di bootstrap si trova nello store per un motivo preciso: essa deve coordinare autenticazione, migrazione, caricamento remoto e aggiornamento della UI. Se questa logica fosse dispersa in piu` componenti React, il rischio di inconsistenze aumenterebbe.

\subsection{Layer di servizi}

Il frontend usa un layer di servizi separato:

\begin{itemize}
  \item \texttt{taskService} applica la logica CRUD dei task;
  \item \texttt{groupService} gestisce CRUD gruppi e calcolo dello stato del gruppo;
  \item \texttt{migrationService} governa import e reimport da \texttt{localStorage};
  \item \texttt{analyticsService} calcola statistiche derivate senza accesso di rete;
  \item \texttt{storageService} realizza persistenza locale e metadati di migrazione.
\end{itemize}

La logica e` quindi stratificata: i componenti React invocano azioni dello store, lo store richiama i servizi di dominio, e solo il \texttt{firestoreService} conosce i dettagli della persistenza remota.

\subsection{Modelli e stato derivato}

I modelli \texttt{Task.js} e \texttt{Group.js} normalizzano la forma dei dati. Alcune regole di business rilevanti sono:

\begin{itemize}
  \item i task ricevono sempre un identificatore, un ordine e una data pianificata;
  \item la priorita` e` vincolata ai valori \texttt{low}, \texttt{medium}, \texttt{high};
  \item i gruppi possiedono milestone con identificatore stabile;
  \item lo stato del gruppo non e` persisito come verita` assoluta, ma ricalcolato dai task del giorno.
\end{itemize}

Quest'ultimo punto e` particolarmente interessante: \texttt{groupService.attachStatuses()} evita ridondanza tra dati persistiti e dati derivati. Un gruppo risulta \texttt{inactive}, \texttt{in\_progress} oppure \texttt{completed} in base ai task associati alla data odierna.

\section{Flussi applicativi principali}

\subsection{Primo accesso e migrazione iniziale}

Il primo login costituisce il flusso di rete piu` ricco. La sequenza implementata e` la seguente:

\begin{enumerate}
  \item l'utente preme ``Continua con Google'';
  \item Firebase Auth autentica l'utente e restituisce il suo profilo;
  \item \texttt{upsertUserProfile()} crea o aggiorna \texttt{users/\{uid\}};
  \item \texttt{migrationService.migrateOnFirstLogin()} verifica se esiste un backup locale importabile;
  \item se la migrazione non e` ancora stata eseguita, task, gruppi e milestone vengono copiati in Firestore tramite batch write;
  \item il profilo utente viene marcato con \texttt{migration.localStorageImported = true};
  \item l'app decide se avviare la UI da cache locale o da stato remoto;
  \item al termine la dashboard viene renderizzata.
\end{enumerate}

La migrazione e` in modalita` \emph{merge only}: i documenti gia` esistenti vengono aggiornati per ID, mentre quelli remoti assenti nel backup locale non vengono cancellati. Questa scelta e` prudente e riduce il rischio di perdita dati.

\subsection{Creazione e aggiornamento dei task}

Quando l'utente crea un task, il percorso logico e`:

\begin{enumerate}
  \item il componente di dashboard invia i dati allo store;
  \item \texttt{taskService.createTask()} completa i campi tecnici mancanti;
  \item il servizio Firestore salva il documento remoto e incrementa la revisione del workspace;
  \item lo store aggiorna lo stato locale;
  \item \texttt{syncLocalBackup()} allinea anche il backup in \texttt{localStorage}.
\end{enumerate}

Un meccanismo analogo viene applicato per update, delete e toggle del completamento. Il vantaggio e` che l'app mantiene tre livelli coerenti: UI React, cache locale del browser e stato persistito su Firestore.

\subsection{Gestione dei gruppi}

La gestione dei gruppi e` leggermente piu` complessa dei task per la presenza delle milestone. In particolare:

\begin{itemize}
  \item la creazione di un gruppo salva sia il documento del gruppo sia la sottocollezione \texttt{milestones};
  \item l'aggiornamento usa \texttt{syncMilestonesForGroup()} per creare, aggiornare o rimuovere milestone obsolete;
  \item la cancellazione di un gruppo implica anche la rimozione dei task collegati e delle milestone.
\end{itemize}

Il flusso di delete e` quindi composto da piu` operazioni coordinate, e lo store si occupa di mantenere consistente il gruppo selezionato nella UI dopo la rimozione.

\subsection{Reimport manuale}

L'app espone anche una funzione di reimport manuale dalla schermata impostazioni. Questo flusso e` utile per riallineare vecchi dati locali con Firestore. Un vincolo importante e` la protezione contro account diversi: se il backup locale risulta associato a un altro \texttt{uid}, il reimport viene bloccato esplicitamente.

\section{Aspetti di coerenza, prestazioni e limiti}

La soluzione adottata presenta diversi punti di forza:

\begin{itemize}
  \item autenticazione delegata a un provider affidabile;
  \item isolamento dei dati tramite regole Firestore per utente;
  \item uso di batch write per operazioni atomiche;
  \item cache locale utile per velocizzare il bootstrap;
  \item separazione abbastanza chiara tra UI, store, servizi di dominio e persistenza.
\end{itemize}

Esistono tuttavia anche alcuni limiti strutturali:

\begin{itemize}
  \item il caricamento dati usa letture esplicite, non listener realtime Firestore;
  \item l'assenza di backend custom rende piu` difficile centralizzare logica avanzata, auditing o elaborazioni server-side;
  \item alcune operazioni complesse, come la cancellazione coordinata di gruppo e task, rimangono gestite interamente nel client;
  \item la sincronizzazione multi-dispositivo si basa sul refresh e sulla persistenza, non su push continuo verso la UI.
\end{itemize}

Dal punto di vista di un corso di reti di calcolatori, l'app rappresenta quindi un buon esempio di architettura cloud \emph{serverless client-driven}: semplice da distribuire, ma meno controllabile di una soluzione con backend applicativo dedicato.

\section{Conclusione}

La parte di rete di Focusly e` costruita attorno a una comunicazione diretta tra frontend React e servizi Firebase. Firebase Authentication gestisce l'identita` dell'utente, mentre Firestore conserva in cloud l'intero workspace personale. Lo store Zustand svolge il ruolo di coordinatore della business logic, orchestrando bootstrap, migrazione, CRUD, sincronizzazione locale e aggiornamento della UI.

In sintesi, l'app non implementa una rete complessa con piu` nodi applicativi proprietari, ma una soluzione moderna basata su servizi cloud gestiti. Il risultato e` un'architettura pulita e adatta a una prima versione del prodotto, nella quale l'aspetto piu` interessante non e` la presenza di un backend tradizionale, bensi` il modo in cui il client organizza autenticazione, persistenza remota e consistenza dei dati tra browser e cloud.

\end{document}
